\subsection{Módulo 2: Reportes y Estadísticas}

\subsubsection{Descripción General}

El módulo de reportes permite visualizar el historial completo de detecciones vehiculares con filtrados avanzados, gráficos estadísticos interactivos y exportación a Excel/CSV.

\subsubsection{Vista Principal - modulo2.blade.php}

\textbf{Archivo}: \texttt{resources/views/modulo2.blade.php} (341 líneas)

\paragraph{Estructura HTML (Líneas 1-180)}

\textbf{Filtros del formulario (L.24-65)}:
\begin{itemize}
    \item \textbf{L.27-30}: Select de cámara con opciones CAM\_001 y CAM\_002
    \item \textbf{L.35}: Input type='date' para fecha\_inicio, valor default: \texttt{date('Y-m-01')}
    \item \textbf{L.40}: Input type='date' para fecha\_fin, valor default: \texttt{date('Y-m-d')}
    \item \textbf{L.45-50}: Select de tipo con opciones: Todos, Auto, Moto, Bus
    \item \textbf{L.55-59}: Select de agrupación con modos: Automático, Por Hora, Por Día
\end{itemize}

\textbf{Canvas para gráficos (L.70-91)}:
\begin{itemize}
    \item \textbf{L.77}: Canvas \texttt{\#dailyChart} para gráfico de barras (300px alto)
    \item \textbf{L.89}: Canvas \texttt{\#typeDistributionChart} para gráfico de torta
\end{itemize}

\textbf{Tabla de registros (L.107-175)}:
\begin{lstlisting}[style=phpstyle, caption={Loop de Detecciones - modulo2.blade.php}]
@forelse($registros ?? [] as $registro)
    @php
        // L.120-126: Mapeo de colores por tipo
        $colorClass = match(strtolower($registro['tipo'])) {
            'auto' => 'bg-blue-500',
            'moto' => 'bg-red-500',
            'bus' => 'bg-green-500',
            default => 'bg-gray-500'
        };
        
        // L.127-128: Determinación de color según confianza
        $conf = floatval($registro['confianza']);
        $confColor = $conf > 90 ? 'bg-green-100' : 
                    ($conf > 75 ? 'bg-yellow-100' : 'bg-red-100');
        
        // L.131-142: Mapeo de colores de vehículo a CSS
        $vehicleColorMap = [
            'rojo' => 'bg-red-500',
            'azul' => 'bg-blue-500',
            'blanco' => 'bg-gray-100 border',
            // ... otros colores
        ];
    @endphp
    
    <tr class="hover:bg-gray-50">
        <td>{{ $registro['fecha'] }}</td>
        <td><span class="{{ $colorClass }}"></span> {{ $registro['tipo'] }}</td>
        <td><span class="{{ $colorBadge }}"></span> {{ $registro['color'] }}</td>
    </tr>
@empty
    <tr><td colspan="6">No se encontraron registros</td></tr>
@endforelse
\end{lstlisting}

\paragraph{JavaScript de Gráficos (L.186-340)}

\textbf{Variables globales (L.188-193)}:
\begin{lstlisting}[language=JavaScript]
let chartBar = null;    // Gráfico de barras
let chartPie = null;    // Gráfico de torta
const getSelectedCamera = () => {
    const params = new URLSearchParams(window.location.search);
    return params.get('camera_id') || 'CAM_002';
};
\end{lstlisting}

\textbf{Función processAndInitCharts() (L.209-333)}:

\textbf{Determinación de modo de visualización (L.211-231)}:
\begin{lstlisting}[language=JavaScript]
const viewMode = urlParams.get('view_mode') || 'auto';
let isMultiDay = false;

if (viewMode === 'daily') {
    isMultiDay = true;  // Agrupar por día
} else if (viewMode === 'hourly') {
    isMultiDay = false;  // Agrupar por hora 0-23
} else {
    // Modo Automático: detecta si hay múltiples días
    const uniqueDates = new Set(data.map(item => item.fecha.split(' ')[0]));
    if (uniqueDates.size > 1) isMultiDay = true;
}
\end{lstlisting}

\textbf{Procesamiento de datos (L.233-261)}:
\begin{lstlisting}[language=JavaScript]
const stats = {};
const typesCount = { 'Auto': 0, 'Moto': 0, 'Bus': 0, 'Otro': 0 };

data.forEach(item => {
    const [datePart, timePart] = item.fecha.split(' ');
    const hour = timePart.split(':')[0];
    
    // Clave de agrupación: fecha completa O hora
    const key = isMultiDay ? datePart : `${hour}:00`;
    
    // Inicializar contador si no existe
    if (!stats[key]) {
        stats[key] = { 'Auto': 0, 'Moto': 0, 'Bus': 0, 'Otro': 0 };
    }
    
    // Incrementar contador
    stats[key][item.tipo]++;
    typesCount[item.tipo]++;
});
\end{lstlisting}

\textbf{Creación de gráfico de barras (L.292-312)}:
\begin{lstlisting}[language=JavaScript]
chartBar = new Chart(ctxBar, {
    type: 'bar',
    data: {
        labels: sortedKeys,  // Ordenadas cronológicamente
        datasets: [
            { label: 'Autos', data: dataAuto, backgroundColor: '#3b82f6' },
            { label: 'Motos', data: dataMoto, backgroundColor: '#ef4444' },
            { label: 'Buses', data: dataBus, backgroundColor: '#10b981' }
        ]
    },
    options: {
        responsive: true,
        maintainAspectRatio: false,
        scales: {
            y: { beginAtZero: true, ticks: { stepSize: 1 } }
        }
    }
});
\end{lstlisting}

\subsubsection{Controlador ReporteController}

\textbf{Archivo}: \texttt{app/Http/Controllers/ReporteController.php} (171 líneas)

\paragraph{Método obtenerRegistrosFiltrados() (L.12-73)}

\textbf{Carga de XML (L.15-19)}:
\begin{lstlisting}[style=phpstyle]
$xmlPath = storage_path('app/vehiculos_db.xml');
if (file_exists($xmlPath)) {
    $xmlContent = file_get_contents($xmlPath);
    $xml = simplexml_load_string($xmlContent);
}
\end{lstlisting}

\textbf{Construcción de registro (L.23-30)}:
\begin{lstlisting}[style=phpstyle]
$registro = [
    'fecha' => (string)$det->fecha,
    'tipo'  => (string)$det->tipo,
    'confianza' => (float)$det->confianza,
    'color' => isset($det->color) ? (string)$det->color : 'desconocido',
    'camara' => isset($det->camara) ? (string)$det->camara : 'CAM_002',
    'nombre_camara' => isset($det->nombre_camara) ? 
        (string)$det->nombre_camara : 'Skyline Cochabamba'
];
\end{lstlisting}

\textbf{Aplicación de filtros (L.34-60)}:
\begin{itemize}
    \item \textbf{L.34-38}: Filtro por camera\_id - compara string exacto
    \item \textbf{L.40-44}: Filtro por tipo - compara string exacto
    \item \textbf{L.46-52}: Filtro por fecha inicio:
    \begin{lstlisting}[style=phpstyle]
$fechaRegistro = strtotime($registro['fecha']);
$fechaInicio = strtotime($request->fecha_inicio . ' 00:00:00');
if ($fechaRegistro < $fechaInicio) {
    $incluir = false;
}
    \end{lstlisting}
    \item \textbf{L.54-60}: Similar para fecha\_fin, añade '23:59:59'
\end{itemize}

\paragraph{Método exportarExcel() (L.82-140)}

\textbf{Setup de archivo (L.84-92)}:
\begin{lstlisting}[style=phpstyle]
$registros = $this->obtenerRegistrosFiltrados($request);
$nombreArchivo = 'reportes_' . date('Y-m-d_His') . '.csv';
$rutaTemporal = storage_path('app/temp');

if (!is_dir($rutaTemporal)) {
    mkdir($rutaTemporal, 0755, true);
}

$rutaCompleta = $rutaTemporal . '/' . $nombreArchivo;
\end{lstlisting}

\textbf{Escritura CSV con BOM UTF-8 (L.100-122)}:
\begin{lstlisting}[style=phpstyle]
$archivo = fopen($rutaCompleta, 'w');

// BOM para Excel UTF-8
fprintf($archivo, chr(0xEF).chr(0xBB).chr(0xBF));

// Encabezados
fputcsv($archivo, [
    'Fecha/Hora', 'Tipo Vehículo', 'Confianza (%)', 
    'Color', 'Cámara', 'Nombre Cámara'
], ',');

// Datos
foreach ($registros as $registro) {
    fputcsv($archivo, [
        $registro['fecha'],
        $registro['tipo'],
        number_format($registro['confianza'], 2),
        $registro['color'],
        $registro['camara'],
        $registro['nombre_camara']
    ], ',');
}

fclose($archivo);
\end{lstlisting}

\textbf{Descarga y auto-eliminación (L.127)}:
\begin{lstlisting}[style=phpstyle]
return response()->download($rutaCompleta, $nombreArchivo)
    ->deleteFileAfterSend(true);
\end{lstlisting}

\newpage
